\tikzset{
    mz base/.style={
        scale=1.3, font=\normalsize, >={Stealth[length=2mm]},
        photon/.style={red!70!black},
        port/.style={photon, thick, dashed, ->},
        upper/.style={blue!70!black, thick, dashed, ->},
        lower/.style={orange!85!black, thick, dashed, ->},
        upper label/.style={note, align=right, text=blue!70!black},
        lower label/.style={note, text=orange!85!black},
        detector/.style={thick, draw=black!80, fill=gray!20},
        note/.style={font=\small},
        mirror face/.style={line width=1.6pt, black!80},
    },
    % highlighted route: solid, thicker, with a soft coloured glow
    mz on/.style={solid, line width=1.4pt,
    preaction={draw=#1, -, line width=6pt, opacity=0.25}},
    % route not taken: thin grey dashes
    mz off/.style={draw=black!25, thin, dashed},
    mz route 1/.style={
        upper/.append style={mz on=blue!50},
        lower/.append style={mz off},
        lower label/.append style={text=black!35},
        port/.append style={blue!70!black, mz on=blue!50},
    },
    mz route 2/.style={
        lower/.append style={mz on=orange!60},
        upper/.append style={mz off},
        upper label/.append style={text=black!35},
        port/.append style={orange!85!black, mz on=orange!60},
    },
    pics/bs/.style={code={
            \draw[thick, fill=cyan!10] (-0.35,-0.35) rectangle (0.35,0.35);
            \draw[very thick, cyan!50!black] (-0.35,-0.35) -- (0.35,0.35);
    }},
    pics/mirror back up/.style={code={
            \draw[mirror face] (-0.38,-0.38) -- (0.38,0.38);
            \foreach \t in {-0.28,-0.12,0.04,0.20,0.36}
            \draw[black!55] (\t,\t) -- ++(-0.12,0.12);
    }},
    pics/mirror back down/.style={code={
            \draw[mirror face] (-0.38,-0.38) -- (0.38,0.38);
            \foreach \t in {-0.36,-0.20,-0.04,0.12,0.28}
            \draw[black!55] (\t,\t) -- ++(0.12,-0.12);
    }},
}
\newcommand{\mzbody}{%
    \coordinate (bs1) at (2.9,0);
    \coordinate (m1)  at (2.9,2.6);
    \coordinate (m2)  at (7.0,0);
    \coordinate (bs2) at (7.0,2.6);
    \node[draw, thick, fill=yellow!20, align=center, font=\small,
    minimum height=8mm, inner xsep=2pt] (src) at (0.7,0) {single-photon\\source};
    \draw[photon, line width=1.4pt, ->] (src.east) -- ($(bs1)+(-0.37,0)$);
    \draw[photon, thick, smooth, samples=80, domain=-0.35:0.35, shift={(1.95,0)}]
    plot (\x, {0.18*exp(-16*\x*\x)*sin(2100*\x)});
    \node[photon, above] at (1.95,0.2) {$h\nu$};
    \draw[upper] ($(bs1)+(0,0.35)$) -- ($(m1)+(0,-0.12)$)
    node[midway, left, upper label] {upper arm\\(reflected)};
    \draw[upper] ($(m1)+(0.12,0)$) -- ($(bs2)+(-0.37,0)$);
    \draw[lower] ($(bs1)+(0.35,0)$) -- ($(m2)+(-0.12,0)$)
    node[pos=0.45, below, lower label] {lower arm (transmitted)};
    \draw[lower] ($(m2)+(0,0.12)$) -- ($(bs2)+(0,-0.37)$);
    \pic at (bs1) {bs};
    \pic at (bs2) {bs};
    \pic at (m1) {mirror back up};
    \pic at (m2) {mirror back down};
    \node[below=1mm, note] at ($(bs1)+(0,-0.35)$) {$BS_1$};
    \node[above left, note] at ($(bs2)+(-0.3,0.3)$) {$BS_2$};
    \node[above left, note] at ($(m1)+(-0.1,0.1)$) {$M_1$};
    \node[below right, note] at ($(m2)+(0.1,-0.1)$) {$M_2$};
    \draw[detector] (8.6,2.1) arc[start angle=-90, end angle=90, radius=0.5] -- cycle;
    \node at (8.84,2.6) {$D_1$};
    \draw[detector] (6.5,4.0) arc[start angle=180, end angle=0, radius=0.5] -- cycle;
    \node at (7.0,4.2) {$D_2$};
    \draw[port] ($(bs2)+(0.35,0)$) -- (8.58,2.6);
    \draw[port] ($(bs2)+(0,0.35)$) -- (7.0,3.98);
}

\subsubsection{Construction of the Mach-Zehnder interferometer}\label{sec:from-wave-particle-duality-to-quantum-states:beam-splitters-and-the-mach-zehnder-interferometer:construction-of-the-mach-zehnder-interferometer}

We have just discovered that a single photon entering a 50/50 beam splitter behaves, statistically, like a fair coin toss:
\begin{equation*}
    P(D_1)=P(D_2)=\frac{1}{2}
\end{equation*}
But we still don't know whether this randomness is equivalent to classical ignorance.
\textbf{\emph{How can we test whether the coin-toss analogy is sufficient?}}
We can construct a new experiment using \hl{two beam splitters instead of one.}
This arrangement is called the \definition{Mach-Zehnder Interferometer}.

\highspace
\textcolor{Green3}{\faIcon{tools} \textbf{The experimental setup.}} The \definition{Mach-Zehnder Interferometer} is an experimental apparatus used to study \textbf{interference between two alternative paths} that light (or a quantum particle) can take. It consists of two 50/50 beam splitters, two ordinary mirrors, a single-photon source, and two detectors.

\begin{figure}[htbp]
    \centering
    \tikzsetnextfilename{mach-zehnder-interferometer}
    \begin{tikzpicture}[
            scale=1.3,
            font=\normalsize,
            >={Stealth[length=2mm]},
            photon/.style={red!70!black},
            port/.style={photon, thick, dashed, ->},
            upper/.style={blue!70!black, thick, dashed, ->},
            lower/.style={orange!85!black, thick, dashed, ->},
            detector/.style={thick, draw=black!80, fill=gray!20},
            note/.style={font=\small},
            % 50/50 beam splitter: cube with a diagonal, centred on the given coordinate
            pics/bs/.style={code={
                    \draw[thick, fill=cyan!10] (-0.35,-0.35) rectangle (0.35,0.35);
                    \draw[very thick, cyan!50!black] (-0.35,-0.35) -- (0.35,0.35);
            }},
            % ideal mirror: thick reflecting face along "/" with hatching on the back side
            mirror face/.style={line width=1.6pt, black!80},
            pics/mirror back up/.style={code={
                    \draw[mirror face] (-0.38,-0.38) -- (0.38,0.38);
                    \foreach \t in {-0.28,-0.12,0.04,0.20,0.36}
                    \draw[black!55] (\t,\t) -- ++(-0.12,0.12);
            }},
            pics/mirror back down/.style={code={
                    \draw[mirror face] (-0.38,-0.38) -- (0.38,0.38);
                    \foreach \t in {-0.36,-0.20,-0.04,0.12,0.28}
                    \draw[black!55] (\t,\t) -- ++(0.12,-0.12);
            }},
        ]
        % geometry: BS1 bottom left, M1 top left, M2 bottom right, BS2 top right
        \coordinate (bs1) at (2.9,0);
        \coordinate (m1)  at (2.9,2.6);
        \coordinate (m2)  at (7.0,0);
        \coordinate (bs2) at (7.0,2.6);

        % source and input beam with one wave packet
        \node[draw, thick, fill=yellow!20, align=center, font=\small,
        minimum height=8mm, inner xsep=2pt] (src) at (0.7,0) {single-photon\\source};
        \draw[photon, thick, ->] (src.east) -- ($(bs1)+(-0.37,0)$);
        \draw[photon, thick, smooth, samples=80, domain=-0.35:0.35, shift={(1.95,0)}]
        plot (\x, {0.18*exp(-16*\x*\x)*sin(2100*\x)});
        \node[photon, above] at (1.95,0.2) {$h\nu$};

        % the two available routes (not two pieces of a photon)
        \draw[upper] ($(bs1)+(0,0.35)$) -- ($(m1)+(0,-0.12)$)
        node[midway, left, note, align=right, text=blue!70!black] {upper arm\\(reflected)};
        \draw[upper] ($(m1)+(0.12,0)$) -- ($(bs2)+(-0.37,0)$);
        \draw[lower] ($(bs1)+(0.35,0)$) -- ($(m2)+(-0.12,0)$)
        node[pos=0.45, below, note, text=orange!85!black] {lower arm (transmitted)};
        \draw[lower] ($(m2)+(0,0.12)$) -- ($(bs2)+(0,-0.37)$);

        % optical elements
        \pic at (bs1) {bs};
        \pic at (bs2) {bs};
        \pic at (m1) {mirror back up};
        \pic at (m2) {mirror back down};
        \node[below=1mm, note] at ($(bs1)+(0,-0.35)$) {$BS_1$};
        \node[above left, note] at ($(bs2)+(-0.3,0.3)$) {$BS_2$};
        \node[above left, note] at ($(m1)+(-0.1,0.1)$) {$M_1$};
        \node[below right, note] at ($(m2)+(0.1,-0.1)$) {$M_2$};

        % detectors (D-shaped, flat face towards the beam) and output ports
        \draw[detector] (8.6,2.1) arc[start angle=-90, end angle=90, radius=0.5] -- cycle;
        \node at (8.84,2.6) {$D_1$};
        \draw[detector] (6.5,4.0) arc[start angle=180, end angle=0, radius=0.5] -- cycle;
        \node at (7.0,4.2) {$D_2$};
        \draw[port] ($(bs2)+(0.35,0)$) -- (8.58,2.6);
        \draw[port] ($(bs2)+(0,0.35)$) -- (7.0,3.98);
    \end{tikzpicture}
    \caption{\textbf{Mach-Zehnder interferometer.} $BS_1$ and $BS_2$ are 50/50
        beam splitters; $M_1$ and $M_2$ are fully reflecting mirrors. The colored lines
        show the two available routes, \emph{not} two pieces of a photon; the red dashed lines
    mark the two output ports leading to the detectors $D_1$ and $D_2$.}
    \label{fig:mach-zehnder-interferometer}
\end{figure}

\noindent
Let's analyze the Mach-Zehnder interferometer in detail:
\begin{itemize}
    \item \important{The first beam splitter ${BS}_1$ creates two possible paths for the photon.} A single photon reaches the first beam splitter, $BS_1$. Just as in the previous section, there are two possible outgoing directions:
        \begin{itemize}
            \item \textbf{Upper arm}: the reflected output.
            \item \textbf{Lower arm}: the transmitted output.
        \end{itemize}
        We know that, if these paths were separately measured immediately after ${BS}_1$, each would have a 50\% detection probability. But in this experiment, we \textbf{do not place detectors on the two arms}, instead we allow both possible routes to continue.

        \textcolor{DarkOrange3}{\faIcon{exclamation-triangle} \textbf{Remark:}} drawing two paths does not mean the photon has been divided into two half-photons. Nor have we established that it secretly chose one definite path. For now, we simply identify the two alternatives available in the apparatus.

    \item \important{The mirrors $M_1$ and $M_2$ redirect the two paths.} Next, each path encounters an ordinary mirror:
        \begin{itemize}
            \item $M_1$ redirects the upper arm toward the second beam splitter.
            \item $M_2$ redirects the lower arm toward the second beam splitter.
        \end{itemize}
        We assume the mirrors are ideal, with 100\% reflectivity. Unlike a beam splitter, \textbf{an ordinary mirror does not create two alternative outgoing paths} in this setup; it redirects the incident beam. The goal of the mirrors is simply to bring the two arms together at the second beam splitter, $BS_2$.

    \item \important{The second beam splitter ${BS}_2$ brings the alternatives together.} This is the essential new ingredient, because in our previous experiment, we had:
        \begin{equation*}
            \text{Photon}\rightarrow BS_1\rightarrow
            \begin{cases}
                D_1\\
                D_2
            \end{cases}
        \end{equation*}
        But now, the photon has two alternative paths that are recombined at $BS_2$:
        \begin{figure}[htbp]
            \centering
            \tikzsetnextfilename{mach-zehnder-paths}
            \begin{tikzpicture}[
                    >={Stealth[length=2mm]},
                    node distance=9mm and 6mm,
                    box/.style={draw, thick, rounded corners=2pt, minimum height=7mm,
                    minimum width=16mm, align=center, fill=#1},
                    box/.default=white,
                    bs/.style={box=cyan!10, minimum width=12mm},
                    det/.style={box=gray!20, minimum width=10mm},
                    flow/.style={line width=1pt, ->},
                    upper/.style={flow, blue!70!black},
                    lower/.style={flow, orange!85!black},
                ]
                % top to bottom: photon, BS1, the two arms, BS2, the two detectors
                \node[box=yellow!20]              (ph)  {Photon};
                \node[bs, below=of ph]            (bs1) {$\mathrm{BS}_1$};
                \node[box=blue!8, draw=blue!70!black,
                below left=of bs1]          (up)  {Upper arm};
                \node[box=orange!10, draw=orange!85!black,
                below right=of bs1]         (lo)  {Lower arm};
                \node[bs, below=of bs1, yshift=-16mm] (bs2) {$\mathrm{BS}_2$};
                \node[det, below left=of bs2]     (d1)  {$D_1$};
                \node[det, below right=of bs2]    (d2)  {$D_2$};

                \draw[flow, red!70!black] (ph) -- (bs1);
                \draw[upper] (bs1) -- (up);
                \draw[lower] (bs1) -- (lo);
                \draw[upper] (up)  -- (bs2);
                \draw[lower] (lo)  -- (bs2);
                \draw[flow]  (bs2) -- (d1);
                \draw[flow]  (bs2) -- (d2);
            \end{tikzpicture}
        \end{figure}

        Both routes lead to the \textbf{same second beam splitter}, $BS_2$, and after this point, we place two detectors, $D_1$ and $D_2$, to record the outcome. This means that, for a detection at one of the final outputs, the apparatus contains \textbf{two possible contributing routes} from the original source:
        \begin{itemize}
            \item \textbf{Route 1:} $S\rightarrow BS_1\rightarrow M_1\rightarrow BS_2\rightarrow D_1$ or $D_2$.
                \begin{figure}[H]
                    \centering
                    \tikzsetnextfilename{mach-zehnder-route-1}
                    \begin{tikzpicture}[mz base, mz route 1]
                        \mzbody
                    \end{tikzpicture}
                \end{figure}
            \item \textbf{Route 2:} $S\rightarrow BS_1\rightarrow M_2\rightarrow BS_2\rightarrow D_1$ or $D_2$.
                \begin{figure}[H]
                    \centering
                    \tikzsetnextfilename{mach-zehnder-route-2}
                    \begin{tikzpicture}[mz base, mz route 2]
                        \mzbody
                    \end{tikzpicture}
                \end{figure}
        \end{itemize}
\end{itemize}
\textcolor{Green3}{\faIcon{question-circle} \textbf{What are we trying to discover?}} We now have two beam splitters, each individually characterized as 50/50. If we treat them as ordinary classical devices that make independent random choices, \hl{what should the final detectors observe?} Would we still expect:
\begin{equation*}
    P(D_1)=P(D_2)=\frac{1}{2} \;?
\end{equation*}
Or could \hl{bringing the two paths together change the outcome statistics?}

\newpage

\begin{flushleft}
    \textcolor{Green3}{\faIcon{times-circle} \textbf{Why two independent 50/50 coin tosses give the wrong prediction}}
\end{flushleft}
We have constructed the Mach-Zehnder interferometer using two 50/50 beam splitters, two mirrors, and two detectors. Now we want to answer a simple question: \textbf{\emph{what should we observe if each beam splitter acts like a fair coin toss?}}
\begin{itemize}
    \item \important{First beam splitter: the first coin toss.} A single photon enters ${BS}_1$. Under our classical random-choice hypothesis, the photon randomly chooses one of the two available paths:
        \begin{equation*}
            P(U)=\frac{1}{2},\qquad P(L)=\frac{1}{2}
        \end{equation*}
        Where $U$ and $L$ denote the upper and lower arms, respectively. So, according to this model, we imagine that the photon takes \textbf{one definite path}, although we do not know which one. This is exactly the coin-toss analogy: heads (upper arm) or tails (lower arm). But remember that this is our \textbf{hypothesis}, not a proven fact.

    \item \important{Second beam splitter: the second coin toss.} The mirrors redirect the paths toward ${BS}_2$. Let's assume the second beam splitter works like another independent 50/50 coin toss. Regardless of which arm the photon came from, it has:
        \begin{equation*}
            P\left(D_{1} \mid U\right) = P\left(D_{2} \mid U\right) = \frac{1}{2}, \quad P\left(D_{1} \mid L\right) = P\left(D_{2} \mid L\right) = \frac{1}{2}
        \end{equation*}
        Where $P\left(D_{1} \mid U\right)$ means \hl{the probability of reaching detector $D_1$ given that the photon followed the upper arm.}
\end{itemize}
\textcolor{Green3}{\faIcon{tools} \textbf{Calculate the four possible outcomes.}} If the two coin tosses are independent, the probability of a particular sequence is the product of the two probabilities:

\begin{figure}[H]
    \centering
    \tikzsetnextfilename{classical-coin-toss-tree}
    \begin{tikzpicture}[
            >={Stealth[length=2mm]},
            font=\small,
            upcol/.style={blue!70!black},
            lowcol/.style={orange!85!black},
            branch/.style={thick, ->},
            box/.style={draw, thick, rounded corners=2pt, minimum width=10mm,
            minimum height=6mm, inner sep=2pt},
            bs/.style={box, fill=cyan!10, draw=cyan!50!black},
            arm U/.style={box, upcol, fill=blue!8},
            arm L/.style={box, lowcol, fill=orange!10},
            det1/.style={box, draw=black!80, fill=green!12, minimum width=8mm},
            det2/.style={box, draw=black!80, fill=violet!12, minimum width=8mm},
            prob/.style={inner sep=1.5pt},
            leaf/.style={anchor=west, inner sep=2pt},
            head/.style={font=\footnotesize\bfseries, text=black!60},
        ]
        % --- column headers (over the two tosses and the leaf column) ------
        \node[head] at (1.25,3.1) {toss at $\mathrm{BS}_1$};
        \node[head] at (3.75,3.1) {toss at $\mathrm{BS}_2$};
        \node[head, anchor=west] at (5.45,3.1) {probability of the route};

        % --- tree nodes ----------------------------------------------------
        \node[bs]    (bs1) at (0,0)      {$\mathrm{BS}_1$};
        \node[arm U] (U)   at (2.5,1.5)  {$U$};
        \node[arm L] (L)   at (2.5,-1.5) {$L$};
        \node[det1]  (UD1) at (5,2.25)   {$D_1$};
        \node[det2]  (UD2) at (5,0.75)   {$D_2$};
        \node[det1]  (LD1) at (5,-0.75)  {$D_1$};
        \node[det2]  (LD2) at (5,-2.25)  {$D_2$};

        % --- first toss ----------------------------------------------------
        \draw[branch, upcol]  (bs1.north east) -- node[prob, above left=-1pt]
        {$\tfrac12$ heads} (U.west);
        \draw[branch, lowcol] (bs1.south east) -- node[prob, below left=-1pt]
        {$\tfrac12$ tails} (L.west);

        % --- second toss: independent, 1/2 whichever arm was taken ---------
        \foreach \arm/\col in {U/upcol, L/lowcol}{
            \draw[branch, \col] (\arm.east) -- node[prob, above, pos=0.45]
            {$\tfrac12$} (\arm D1.west);
            \draw[branch, \col] (\arm.east) -- node[prob, below, pos=0.45]
            {$\tfrac12$} (\arm D2.west);
        }

        % --- leaf probabilities (same notation as the notes) ---------------
        \foreach \arm/\d/\col in {U/1/upcol, U/2/upcol, L/1/lowcol, L/2/lowcol}{
            \node[leaf, \col] at (\arm D\d.east)
            {$P(\arm\text{ then }D_{\d})=\tfrac12\cdot\tfrac12=\tfrac14$};
        }

        % --- totals: each detector is reached through both arms ------------
        \draw[black!25] (-0.6,-2.95) -- (10.5,-2.95);
        \node[det1, thick, inner xsep=5pt] at (2.2,-3.75)
        {$P(D_1)=\underset{\textcolor{blue!70!black}{U}}{\textcolor{blue!70!black}{\tfrac14}}
            +\underset{\textcolor{orange!85!black}{L}}{\textcolor{orange!85!black}{\tfrac14}}
        =\tfrac12$};
        \node[det2, thick, inner xsep=5pt] at (7.7,-3.75)
        {$P(D_2)=\underset{\textcolor{blue!70!black}{U}}{\textcolor{blue!70!black}{\tfrac14}}
            +\underset{\textcolor{orange!85!black}{L}}{\textcolor{orange!85!black}{\tfrac14}}
        =\tfrac12$};
    \end{tikzpicture}
\end{figure}

\noindent
Notice that \textbf{each final detector can be reached through either of the two arms}, so we can sum the probabilities for each detector:
\begin{equation*}
    P(D_1) = P(U\text{ then }D_1) + P(L\text{ then }D_1) = \frac{1}{4} + \frac{1}{4} = \frac{1}{2} = 50\%
\end{equation*}
\begin{equation*}
    P(D_2) = P(U\text{ then }D_2) + P(L\text{ then }D_2) = \frac{1}{4} + \frac{1}{4} = \frac{1}{2} = 50\%
\end{equation*}
Even though we have two beam splitters, the final statistics still resemble those of a fair coin toss.

\highspace
\textcolor{Red2}{\faIcon{exclamation-triangle} \textbf{But what happens experimentally?}} In a properly aligned Mach-Zehnder interferometer, with coherent paths and suitable relative phase, we can observe:
\begin{equation*}
    P(D_1)=100\%,
    \qquad
    P(D_2)=0\%
\end{equation*}
\hl{All successfully detected photons arrive at the same output!} Of course, we can also arrange the interferometer so that the opposite detector receives all the photons. The exact output depends on the relative phase accumulated along the two arms and the beam-splitter configuration. But the important fact is that \textbf{the result can be completely different from 50/50}.

\highspace
\begin{flushleft}
    \textcolor{Green3}{\faIcon{question-circle} \textbf{Why does the classical prediction fail?}}
\end{flushleft}
Let's identify exactly what went wrong. Our classical model \textbf{assumed} that:
\begin{enumerate}
    \item The first beam splitter randomly selects one definite path.
    \item The second beam splitter makes another independent random choice.
    \item We can calculate final probabilities simply by adding the probabilities of different routes.
\end{enumerate}
But the experiment shows that this model is insufficient. \hl{The two alternatives cannot always be treated as mutually exclusive classical routes whose probabilities are independently added.}

\highspace
In quantum mechanics, the two paths contribute \emph{probability amplitudes}, and these amplitudes can \textbf{interfere} when the paths recombine. That is the missing ingredient. So, the failure of this simple two-coin-toss model does not, by itself, disprove every possible hidden-variable interpretation. It shows that \hl{we cannot describe this interferometer as two independent classical random decisions.}

\highspace
And the next question is: \emph{\textbf{how can two individually 50/50 beam splitters produce a deterministic output?}} In this case, ``\emph{deterministic}'' means that all of the photons are detected at one output port, and none are detected at the other.

% \newpage

% \begin{flushleft}
%     \textcolor{Green3}{\faIcon{wave-square} \textbf{Interference at the second beam splitter}}
% \end{flushleft}